\scanpage{85}
\deleted{Soient $K^{\cdot\cdot}$ et $K'^{\cdot\cdot}$.}

Finalement, la donnée de $K^{\cdot\cdot}$ équivaut à celle de

$\alpha)$ $k$, Anneau comm. sur $X$

$\beta)$ $L$, $k$-Module plat (voire inversible voire $k$ lui-$=$\ldots)

$\gamma)$ Une immersion de $L$ dans un $k$-module $\Phi$, $u:L\mathbin{-}\Phi$.

Soit de $\unsure{=|\mathrm{m.}}$ $K'^{\cdot\cdot}$ défini par $(k',L',\Phi',u')$. \deleted{\unsure{ou}} La donnée d'un hom $\lambda$ de $K^{\cdot\cdot}$ dans $K'^{\cdot\cdot}$, compatible avec ttes les structures, équivaut à la donnée d'un syst. d'hom

\diagram{\begin{tikzcd}
k \arrow[r,"\deleted{\unsure{\lambda}}\lambda_0"] & k'
\end{tikzcd}}
hom d'Anneaux

\deleted{\diagram{\begin{tikzcd}
L \arrow[r,"\unsure{\lambda_L}"] & L'
\end{tikzcd}}}

\diagram{\begin{tikzcd}
\Phi \arrow[r,"\lambda_\Phi",dash] & \Phi'
\end{tikzcd}}
$\left\{\text{hom $\lambda_0$-}\unsure{\tfrac12\text{ linéaire}|k\text{ linéaire}}\right.$
\deleted{\unsure{ou}}

avec $\lambda_\Phi(u(L))\subset u'(L')$, donc $\lambda_\Phi$
\deleted{\illegible{4}}
induit un hom (évidemment $\lambda_0$-$\frac12$ lin.)
\deleted{\diagram{\begin{tikzcd}
L \arrow[r,"\unsure{\lambda} ",dash] \arrow[d] & L' \arrow[d] \\
\Phi \arrow[r,dash] & \Phi'
\end{tikzcd}}}
\diagram{\begin{tikzcd}
L \arrow[r,"\lambda_L"] & L'
\end{tikzcd}}

Dans le cas $L\simeq k$, $L'\simeq k'$ donnés, on peut oublier $L,L'$ et \unsure{retenir} $\Phi,\Phi'$ avec sections marquées, et la condition sur $\lambda_\Phi$

\margin{84}
